\section{Introduction}

Predicting drug-target interactions for unseen molecules or proteins is a central goal in computational drug discovery. Because truly prospective evaluation is difficult and expensive, many recent studies report held-out drug, target, or blind splits on datasets derived from BindingDB or DAVIS as proxies for inductive generalization \citep{liu2007bindingdb,davis2011kinase,ahmed2024dtilm,luo2024druglamp,ouyang2024pmmr,liu2025spdti,yu2025gsdti,zhao2025coldstartcpi}. These protocols are useful, but they also encourage a strong implicit assumption: if a method performs well on one cold-start split, then it is broadly robust to the kinds of shift encountered in deployment.

That assumption is weaker than it appears. Different split definitions remove different forms of overlap, and they may favor different inductive biases. A retrieval mechanism can help when nearby support examples remain informative, yet become much less useful when the promoted benchmark introduces temporal, provenance, or assay changes that are not captured by synthetic exclusion rules. Likewise, a target-memory mechanism may appear beneficial on one synthetic regime while failing to transfer across datasets or assay families. If these ranking changes occur in practice, benchmark design is not a technical footnote. It becomes part of the scientific claim.

This paper studies that question directly. Rather than proposing another universal architecture, we ask whether the ranking of a compact DTI model panel is preserved across synthetic cold-start regimes and real distribution shifts. Our starting point is a synthetic reference panel on \texttt{BindingDB\_Kd}, where we compare three representative systems: a plain base encoder, a retrieval-augmented model (\texttt{RAICD}), and a functional target-memory model (\texttt{FTM sparse}). Even here, rankings are already heterogeneous: \texttt{RAICD} improves on \texttt{blind\_start}, loses on \texttt{unseen\_drug}, and \texttt{FTM sparse} is strongest on synthetic \texttt{unseen\_target}.

We then contrast this synthetic picture with a real temporal/provenance benchmark built from the BindingDB patent split (\texttt{patent\_temporal}), using matched 3-seed evaluation. The result is a clear ranking reversal. On synthetic \texttt{blind\_start}, retrieval is beneficial; on patent temporal shift, the plain base model is best, outperforming both \texttt{RAICD} and \texttt{FTM sparse}. Paired bootstrap confidence intervals support both parts of this contrast: the synthetic \texttt{blind\_start} \AUPRC{} delta for \texttt{RAICD - base} remains positive, while the patent-temporal deltas for both \texttt{RAICD - base} and \texttt{FTM - base} remain negative.

To move beyond leaderboard-style reporting, we also diagnose the promoted temporal benchmark by year and overlap bucket. This analysis shows that temporal OOD is not simply ``harder cold-start.'' Instead, the gap is concentrated in specific year bands and overlap regimes, which helps explain why a method that looks strong under one split can lose under another. We further include \texttt{BindingDB\_Ki} and \texttt{DAVIS} as supporting real-OOD screens, showing that the synthetic \texttt{unseen\_target} advantage of \texttt{FTM sparse} does not transfer reliably outside the original benchmark. Recent molecular benchmark studies have made analogous arguments about OOD sensitivity in adjacent tasks \citep{ektefaie2024generalizability,shen2025ddiben,delafuente2025inductive}; our contribution is to show the same issue concretely in DTI.

Our contributions are fourfold. First, we construct a 3-seed synthetic reference panel showing that even standard cold-start evaluation is regime-dependent. Second, we establish a 3-seed real-OOD benchmark on \texttt{BindingDB\_patent} and show a ranking reversal relative to synthetic \texttt{blind\_start}. Third, we provide year-band and overlap-bucket analyses that connect the ranking change to measurable shift axes. Fourth, we show through supporting benchmarks that method gains observed on synthetic \texttt{unseen\_target} are not robust enough to serve as a universal claim.
